\subsection{Scheduling Strategies}
\label{sec:scheduling-strategies}

We implement and evaluate six task scheduling strategies for the query engine's thread pool.
The query engine uses a two-level queue architecture:

\begin{itemize}
    \item \textbf{Admission queue} (shared, bounded): receives tasks from external producers (sources, control messages). Provides backpressure---if full, the source blocks until space is available.
    \item \textbf{Internal queue(s)} (unbounded): hold successor tasks produced during pipeline execution. Depending on the strategy, this is either a single shared queue or one queue per worker thread.
\end{itemize}

When a worker thread is woken by its semaphore, it reads tasks in the following priority order:
\begin{enumerate}
    \item \textbf{Own per-thread queue} (successor tasks from pipeline execution)
    \item \textbf{Steal from other threads' queues} (if work stealing is enabled)
    \item \textbf{Admission queue} (new source data)
\end{enumerate}
This priority ensures that in-flight pipeline work is completed before accepting new input, preventing the accumulation of half-processed data.

\subsubsection{Single-Queue Strategy}

\begin{description}
    \item[\texttt{GLOBAL\_QUEUE}.]
    All worker threads share a single lock-free MPMC queue (Folly \texttt{MPMCQueue}) as the internal queue.
    Both admission tasks from sources and successor tasks from pipeline execution are placed into their respective shared queues.
    Workers block on a single shared semaphore that counts tasks across both queues.

    \textit{Advantages:} Perfect load balancing---no worker can idle while tasks are available.
    Simple implementation with no configuration parameters.

    \textit{Disadvantages:} Every enqueue and dequeue operation contends on the same atomic state.
    At high thread counts, cache line bouncing on the shared queue's head/tail pointers becomes the bottleneck.
    No data locality---a successor task is likely processed by a different thread than the one that produced it, requiring the data to be fetched from a remote L2/L3 cache.
\end{description}

\subsubsection{Per-Thread Queue Strategies}

The following strategies assign each worker thread its own unbounded internal queue (Folly \texttt{UMPMCQueue}) and its own semaphore.
The shared bounded admission queue remains global.
When a source submits a task, it is written to the admission queue and a worker thread's semaphore is signaled via round-robin to ensure fair wakeup distribution.
The woken worker checks its own per-thread queue first; if empty, it falls through to the admission queue.
Note that the scheduling strategy's selection policy is \emph{only} applied to internal (successor) tasks, which are placed directly into a target thread's per-thread queue.
Admission tasks always use round-robin semaphore signaling regardless of the strategy.

\begin{description}
    \item[\texttt{PER\_THREAD\_ROUND\_ROBIN}.]
    Successor tasks are assigned to worker queues in rotating order using an atomic counter:
    \[
        \text{target} = \texttt{counter.fetch\_add(1)} \mod N
    \]
    where $N$ is the number of worker threads.
    This provides $O(1)$ dispatch with a single relaxed atomic operation.

    \textit{Advantages:} Even distribution of tasks across all threads.
    Minimal dispatch overhead (one atomic increment).

    \textit{Disadvantages:} No awareness of queue load---a slow thread accumulates a backlog while fast threads idle.
    No data locality---successor tasks are distributed away from the producing thread.

    \item[\texttt{PER\_THREAD\_SMALLEST\_QUEUE}.]
    Successor tasks are assigned to the worker queue with the fewest pending tasks.
    Queue sizes are tracked via per-queue \texttt{atomic<int64\_t>} counters that are incremented on enqueue and decremented on dequeue, using \texttt{memory\_order\_relaxed}.
    The dispatch scans all $N$ counters to find the minimum:
    \[
        \text{target} = \arg\min_{i \in [0,N)} \texttt{approxSize}[i]
    \]

    \textit{Advantages:} Adaptive load balancing---threads with shorter queues receive more work.

    \textit{Disadvantages:} $O(N)$ dispatch cost (linear scan of all queue sizes).
    Approximate sizes may be stale under high concurrency, leading to suboptimal choices.
    Multiple threads may simultaneously select the same ``smallest'' queue, causing temporary imbalance.
    No data locality.

    \item[\texttt{PER\_THREAD\_CHOOSE\_TWO}.]
    Successor tasks are assigned to the shorter of two randomly sampled queues (the ``power of two choices'' technique~\cite{mitzenmacher2001power}):
    \[
        a, b \sim \text{Uniform}(0, N), \quad \text{target} =
        \begin{cases}
            a & \text{if } \texttt{approxSize}[a] \leq \texttt{approxSize}[b] \\
            b & \text{otherwise}
        \end{cases}
    \]
    A thread-local \texttt{std::mt19937} random number generator avoids shared state.

    \textit{Advantages:} $O(1)$ dispatch with only two atomic reads.
    Provably achieves near-optimal load balancing---the maximum queue length is $O(\log \log N)$ compared to $O(\log N)$ for random single-choice and $O(1)$ for full-scan~\cite{mitzenmacher2001power}.

    \textit{Disadvantages:} Non-deterministic assignment.
    Slightly less balanced than \texttt{SMALLEST\_QUEUE} in low-contention scenarios.
    No data locality.

    \item[\texttt{PRODUCER\_LOCAL}.]
    A worker thread that produces a successor task places it on \emph{its own} per-thread queue, bypassing the dispatch strategy entirely.
    Tasks from external sources (admission queue) and non-worker contexts (e.g., the delayed task submitter, which re-submits deferred tasks after a timeout) still use the underlying strategy's dispatch policy.
    In practice, since the vast majority of internal tasks originate from worker threads, the choice of underlying strategy becomes largely irrelevant when producer-local is enabled---our experiments confirm that all three per-thread strategies perform identically (within 1\%) in this mode.

    \textit{Advantages:} Optimal data locality---after processing a tuple buffer through a pipeline stage, the data is hot in the producing thread's L1/L2 cache.
    By keeping the successor task (next pipeline stage) on the same thread, we avoid cache misses from cross-thread data migration.
    Zero dispatch overhead for internal tasks (no atomic operations, no queue selection).

    \textit{Disadvantages:} Can cause load imbalance if one thread produces more successors than others, as all successor work accumulates on a single thread's queue.
    This can be mitigated by combining it with work stealing.

    \item[\texttt{PER\_THREAD\_ADAPTIVE}.]
    An adaptive strategy that dynamically switches between local and distributed dispatch based on the producer's queue pressure.
    When a worker produces a successor task, it checks its own approximate queue size:
    \[
        \text{target} =
        \begin{cases}
            \text{producer} & \text{if } \texttt{approxSize}[\text{producer}] \leq 1 \quad \text{(queue shrinking/empty)} \\
            \text{choose-two}(a, b) & \text{if } \texttt{approxSize}[\text{candidate}] < \texttt{approxSize}[\text{producer}] \\
            \text{producer} & \text{otherwise (no less-loaded thread found)}
        \end{cases}
    \]
    When the producer's queue is empty or has at most one pending task, the successor stays local for cache locality (like \texttt{PRODUCER\_LOCAL}).
    When the queue is growing (size $> 1$), the strategy samples two random queues via choose-two and offloads the task to the less loaded one---but only if that candidate is actually less loaded than the producer's queue.
    If no better candidate is found, the task stays local.

    \textit{Advantages:} Combines the cache locality of \texttt{PRODUCER\_LOCAL} (when the thread is keeping up with its work) with the load balancing of \texttt{CHOOSE\_TWO} (when the thread is falling behind).
    The transition is automatic and requires no tuning---the queue size acts as a natural backpressure signal.
    Zero overhead in the common case (queue empty, task stays local).

    \textit{Disadvantages:} Slightly higher dispatch cost than \texttt{PRODUCER\_LOCAL} due to the queue size check on every dispatch.
    The threshold ($\leq 1$) is a heuristic---a thread may keep tasks local even when offloading would improve overall throughput, or offload too eagerly when the queue is temporarily at size~2.
    Our experiments show that \texttt{ADAPTIVE} performs comparably to \texttt{PRODUCER\_LOCAL} at low thread counts and slightly better under high load, but does not outperform \texttt{GLOBAL\_QUEUE} on workloads with few successor tasks.
\end{description}

\subsubsection{Work Stealing}

All per-thread queue strategies (including \texttt{PRODUCER\_LOCAL}) support optional work stealing.
When a worker's own queue is empty, it attempts to dequeue a task from other workers' queues in round-robin order before falling back to the shared admission queue:

\begin{enumerate}
    \item Try own per-thread queue (non-blocking dequeue)
    \item If empty and work stealing enabled: try queue $(i+1) \mod N$, $(i+2) \mod N$, \ldots
    \item Fall back to admission queue (spin-wait, guaranteed to have a task due to semaphore invariant)
\end{enumerate}

\textit{Rationale:} Work stealing prevents idle threads when load is unevenly distributed, particularly when combined with \texttt{PRODUCER\_LOCAL} dispatch where all successor tasks accumulate on the producing thread.

\textit{Trade-off:} Stealing introduces contention on victim queues and breaks data locality (the stolen task's data is cold in the stealer's cache).
Our experiments show that work stealing generally hurts performance for stateless workloads where load is naturally balanced through the admission queue's round-robin wakeup, but can help in combination with \texttt{PRODUCER\_LOCAL} for workloads with uneven task production.
